\chapter{Máscaras de Zonas Estables SIFT}
\label{ch:masks}
% ==============================================================

% Robustez de compilación: las capturas de referencia por cámara de este
% capítulo son assets de la campaña de campo, no versionados en el
% repositorio de código. Si faltan, se muestra un marcador en vez de abortar
% la compilación de toda la memoria — en cuanto se añadan al repo (o a la
% carpeta de trabajo de LaTeX), las figuras se renderizan solas sin tocar
% este fichero.
\newcommand{\safeimg}[2][width=\textwidth]{%
  \IfFileExists{#2}{\includegraphics[#1]{#2}}{%
    \fbox{\parbox[c][4cm][c]{\dimexpr\linewidth-2\fboxsep-2\fboxrule\relax}{\centering
      \ttfamily\small Imagen no disponible en el repositorio:\\ \url{#2}}}%
  }%
}

\section{Problema: objetos dinámicos en SIFT}

Sin restricciones de zona, SIFT detecta features en \textbf{cualquier región} de la imagen,
incluyendo objetos dinámicos que producen correspondencias erróneas:

\begin{itemize}[leftmargin=1.5em]
  \item \textbf{Gaviotas y pájaros} en vuelo (zona de cielo/horizonte bajo)
  \item \textbf{Olas y espuma} de mar (bordes dinámicos muy contrastados)
  \item \textbf{Personas y bañistas} en la playa
  \item \textbf{Sombras dinámicas} de nubes
  \item \textbf{Hojas de palmeras} agitadas por el viento (ver corrección CAM~3)
\end{itemize}

Resultado: matches entre objetos en posiciones distintas entre tomas → homografía sesgada
→ imágenes ``giradas'' o desplazadas incorrectamente.

\section{Solución: máscaras de regiones estables}

El fichero \texttt{calibration/alignment\_masks.json} define zonas con \textbf{coordenadas
fraccionarias} $(0.0 \ldots 1.0)$ independientes de la resolución real de la imagen:

\begin{equation}
  x_\text{frac} = \frac{x_\text{px}}{W_\text{imagen}}, \quad
  y_\text{frac} = \frac{y_\text{px}}{H_\text{imagen}}
\end{equation}

\begin{codebox}[alignment\_masks.json — estructura]
\begin{minted}{json}
{
  "CAM_1": {
    "regions": [
      {"label": "horizonte_mar",   "x0": 0.00, "y0": 0.39, "x1": 1.00, "y1": 0.48},
      {"label": "arena_baja",      "x0": 0.00, "y0": 0.73, "x1": 0.40, "y1": 0.88},
      {"label": "vegetacion_dcha", "x0": 0.40, "y0": 0.73, "x1": 1.00, "y1": 1.00}
    ]
  }
}
\end{minted}
\end{codebox}

\section{Zonas por cámara}

\begin{table}[H]
\centering
\caption{Configuración de máscaras SIFT por cámara}
\label{tab:masks}
\begin{tabular}{C{1.5cm} L{3.5cm} C{3.5cm} L{4.5cm}}
\toprule
\textbf{Cámara} & \textbf{Zona 1} & \textbf{Zona 2} & \textbf{Justificación} \\
\midrule
CAM~1 & Horizonte mar (y=0.39--0.48) & Arena baja + Vegetación dcha ampliada (x=0.40--1.00, hasta fondo) & Excluye cielo con gaviotas \\
CAM~2 & Horizonte (y=0.34--0.48)    & Vallas+Vegetación izq (x=0.00--0.28, hasta el borde inferior) & Zona estable con vallas y vegetación fija \\
CAM~3 & Horizonte (y=0.45--0.56)    & Kiosco esquina inf. dcha (x=0.87--1.00, y=0.80--0.93) & Estructura fija; excluye las palmeras (x=0.72--0.87) que se mueven con el viento \\
CAM~4 & Horizonte (y=0.38--0.52)    & Edificio izq (x=0.00--0.28, hasta el borde inferior) & Edificio como anclaje vertical \\
CAM~5 & Horizonte (y=0.36--0.50)    & Zona dcha (x=0.68--1.00)                   & Zona derecha más estable \\
CAM~6 & Horizonte (y=0.40--0.54)    & Edificios izq (x=0.00--0.35)               & \\
\bottomrule
\end{tabular}
\end{table}

\section{Visualización de zonas estables por cámara}

Las figuras siguientes muestran, sobre una imagen real de cada cámara, las regiones activas
(coloreadas semitransparentes) donde SIFT extrae keypoints.
\textbf{Verde} = zona 1 (horizonte/mar), \textbf{naranja} = zona 2 (estructura fija secundaria),
\textbf{cian} = zona 3 (cuando existe). Las regiones sin color son ignoradas por SIFT.

\begin{figure}[H]
  \centering
  \safeimg{mask_cam1.jpg}
  \caption{CAM~1 — horizonte del mar ($y=0{,}39$--$0{,}48$) + arena baja izq + vegetación dcha (hasta fondo).
           Zona diseñada para evitar las gaviotas que vuelan en $y < 0{,}38$.}
  \label{fig:mask-cam1-vis}
\end{figure}

\begin{figure}[H]
  \centering
  \safeimg{mask_cam2.jpg}
  \caption{CAM~2 — horizonte del mar ($y=0{,}34$--$0{,}48$) + zona izquierda con vallas y vegetación
           ($x=0{,}00$--$0{,}28$, hasta el borde inferior). Las vallas y masa vegetal ofrecen features
           estables cuando la visibilidad no permite distinguir el Peñón de Ifach.}
  \label{fig:mask-cam2-vis}
\end{figure}

\begin{figure}[H]
  \centering
  \safeimg{mask_cam3.jpg}
  \caption{CAM~3 — horizonte del mar ($y=0{,}45$--$0{,}56$) + kiosco de playa en la esquina
           inferior derecha ($x=0{,}87$--$1{,}00$, $y=0{,}80$--$0{,}93$). Se excluye deliberadamente
           la franja $x=0{,}72$--$0{,}87$: las palmeras que hay ahí agitan sus hojas con el viento
           y generaban correspondencias SIFT inestables (rotaciones espurias de más de
           $100^\circ$ en las pruebas de alineación).}
  \label{fig:mask-cam3-vis}
\end{figure}

\begin{figure}[H]
  \centering
  \safeimg{mask_cam4.jpg}
  \caption{CAM~4 — horizonte del mar ($y=0{,}38$--$0{,}52$) + paseo marítimo/edificios en
           la esquina izquierda ($x=0{,}00$--$0{,}28$, hasta el borde inferior). La infraestructura
           urbana proporciona features angulares muy discriminantes.}
  \label{fig:mask-cam4-vis}
\end{figure}

\begin{figure}[H]
  \centering
  \safeimg{mask_cam5.jpg}
  \caption{CAM~5 — horizonte del mar ($y=0{,}36$--$0{,}50$) + franja derecha estrecha
           ($x=0{,}68$--$1{,}00$). La zona derecha captura vegetación y estructuras fijas
           con mayor estabilidad que la franja izquierda anterior.}
  \label{fig:mask-cam5-vis}
\end{figure}

\begin{figure}[H]
  \centering
  \safeimg{mask_cam6.jpg}
  \caption{CAM~6 — horizonte del mar ($y=0{,}40$--$0{,}54$) + bloque de edificios en la
           esquina izquierda ($x=0{,}00$--$0{,}35$). Los edificios tienen fachadas con
           ventanas y balcones que generan features muy estables.}
  \label{fig:mask-cam6-vis}
\end{figure}

\begin{warnbox}[Nota sobre CAM\_2 — zona \texttt{vallas\_vegetacion}]
La región etiquetada como \texttt{vallas\_vegetacion} ($x=0{,}00$--$0{,}28$, $y=0{,}41$--$1{,}00$)
cubre la franja izquierda donde aparecen las vallas y la vegetación costera, que son estructuras
estáticas presentes en todas las condiciones de visibilidad. Se extiende hasta el borde inferior
de la imagen. En días de alta visibilidad, el Peñón de Ifach puede asomar en esa misma zona como
referencia adicional de gran contraste.
\end{warnbox}

\section{Corrección CAM\_1: problema de gaviotas}

\begin{figure}[H]
  \centering
  \begin{subfigure}[b]{0.48\textwidth}
    \safeimg{img_cam1_mask_actual.jpg}
    \caption{Máscara anterior — zona ``horizonte'' comienza en $y=0{,}24$, alcanzando el cielo con gaviotas}
  \end{subfigure}
  \hfill
  \begin{subfigure}[b]{0.48\textwidth}
    \safeimg{img_cam1_mask_propuesta.jpg}
    \caption{Máscara corregida — \texttt{horizonte\_mar} comienza en $y=0{,}39$ (mar bajo), más arena baja y vegetación hasta el fondo}
  \end{subfigure}
  \caption{Comparativa de máscaras SIFT para CAM~1 (imagen 4608$\times$2682~px)}
  \label{fig:mask-cam1}
\end{figure}

\begin{infobox}[Por qué CAM\_1 es problemática]
La cámara apunta hacia el norte sobre la playa. En imágenes de mañana, 5--10 gaviotas
sobrevuelan habitualmente en la banda $y \approx 0{,}14 \ldots 0{,}30$ (cielo bajo).

SIFT detecta las aves como features de alto contraste con gran confianza. Al estar en
posiciones distintas entre tomas, generan matches cruzados que sesgan la homografía.

La corrección: \texttt{horizonte\_mar} ($y=0{,}39 \ldots 0{,}48$) está en la franja de agua
cerca del horizonte, donde los features son reproducibles y estables.
\end{infobox}

\subsection{Coordenadas en píxeles (imagen 4608$\times$2682)}

\begin{table}[H]
\centering
\caption{Conversión fracción → píxeles para CAM\_1}
\begin{tabular}{L{4cm} C{4cm} C{4cm}}
\toprule
\textbf{Zona} & \textbf{Esquina sup-izq (px)} & \textbf{Esquina inf-der (px)} \\
\midrule
\texttt{horizonte\_mar}   & (0,~1046)   & (4608,~1287) \\
\texttt{arena\_baja}      & (0,~1957)   & (1843,~2360) \\
\texttt{vegetacion\_dcha} & (1843,~1957)& (4608,~2682) \\
\bottomrule
\end{tabular}
\end{table}

\section{Cómo ajustar una máscara}

\begin{enumerate}[leftmargin=1.5em]
  \item Abrir una imagen real de la cámara en un visor que muestre coordenadas de píxel.
  \item Identificar las regiones estables (sin gaviotas, olas ni personas).
  \item Anotar las coordenadas en píxeles de las esquinas.
  \item Convertir a fracciones: $x_f = x_\text{px}/W$, $y_f = y_\text{px}/H$.
  \item Editar \texttt{calibration/alignment\_masks.json}.
  \item Ejecutar el script de verificación:
\end{enumerate}

\begin{codebox}[Script de verificación visual de máscaras]
\begin{minted}{python}
import cv2, json, numpy as np

IMG   = "proces_images/data/camera1/IMAGEN_REF.jpg"
MASKS = "calibration/alignment_masks.json"
CAM   = "CAM_1"

img = cv2.imread(IMG)
h, w = img.shape[:2]
overlay = img.copy()

with open(MASKS) as f:
    cfg = json.load(f)

for r in cfg[CAM]["regions"]:
    x0, y0 = int(r["x0"]*w), int(r["y0"]*h)
    x1, y1 = int(r["x1"]*w), int(r["y1"]*h)
    cv2.rectangle(overlay, (x0,y0), (x1,y1), (0,255,0), -1)
    cv2.putText(overlay, r["label"], (x0+10, y0+30),
                cv2.FONT_HERSHEY_SIMPLEX, 0.8, (0,0,0), 2)

result = cv2.addWeighted(img, 0.5, overlay, 0.5, 0)
cv2.imwrite("/tmp/mask_check.jpg", result)
print(f"Guardado en /tmp/mask_check.jpg  ({w}x{h})")
\end{minted}
\end{codebox}

% ==============================================================
